\section{Related Work}

Overcoming the medical data silo has driven significant research into synthetic generation. Deep Generative models, such as medGAN \cite{choi2017} and Variational Autoencoders (VAEs), are frequently used to synthesize tabular EHRs. While statistically powerful, these purely data-driven models often struggle to maintain logical and physiological constraints (e.g., preventing contradictory laboratory values). In contrast, our approach utilizes a rule-based, expert-driven physiological generation pipeline, supporting clinical coherence across multiple data modalities, making it suitable for validating diagnostic algorithms. However, rule-based engines carry the opposite risk: without empirical grounding, the generated distributions may not reflect real clinical populations. Public Emergency Department resources, such as MIMIC-IV-ED and open-access hospital admission datasets, are increasingly used as reference distributions to quantify this gap, a practice we adopt to validate our generator.

Recent literature in medical AI emphasizes the necessity of multimodal fusion—synthesizing high-frequency sensor data with low-frequency clinical information \cite{baltrusaitis2019}. While most IoMT frameworks focus on localized time-series analysis (e.g., detecting arrhythmias) \cite{rahmani2018}, our synthetic dataset simulates the spatial fusion of these modalities into a unified tabular snapshot, representing the exact moment of emergency triage.

Tabular data remains a challenging domain for Deep Learning, historically dominated by tree-based methods. Gradient-boosted decision trees, most notably XGBoost, remain a strong baseline on structured clinical data due to their native handling of heterogeneous, sparse feature spaces. More recently, architectures like the FT-Transformer \cite{gorishniy2021} have shown that projecting tabular features into embedding tokens allows models to capture complex, non-linear interactions via Self-Attention. Because neither family of models dominates uniformly across tabular benchmarks, we compare both paradigms in this work rather than restricting the evaluation to deep architectures alone. Furthermore, to ensure trust in these models, SHAP \cite{lundberg2017} has become a standard for Explainable AI (XAI), providing local and global feature importance that we use to validate the physiological integrity of our synthetic data across both model families.

A pervasive issue in medical Deep Learning is the class imbalance inherent to real-world epidemiological data, where rare, life-threatening conditions are outnumbered by benign cases. Traditional resampling techniques, such as SMOTE \cite{chawla2002}, often fail in multimodal contexts because they interpolate in the feature space without respecting physiological boundaries, generating clinically impossible synthetic samples. Our approach circumvents this limitation. By leveraging an expert-driven rule-based engine, we directly generate the minority classes (e.g., severe sepsis or hemorrhagic strokes) at scale. Combined with cost-sensitive learning algorithms, such as our weighted Cross-Entropy implementation, this ensures that the neural networks learn accurate, discriminative boundaries for critical conditions.